Let us introduce the last piece of terminology and main object of study in this paper.

\begin{df}
For any $m, N\in\N$ with $1\leq m\leq N$, $\lambda\in\GT_N$, def\/ine
\begin{gather}\label{normalmacdonald}
P_{\lambda}(x_1, \dots, x_m; N, q, t) \myeq \frac{P_{\lambda}\big(x_1, \dots, x_m, 1, t, \dots, t^{N-m-1}; q, t\big)}{P_{\lambda}\big(1, t, t^2, \dots, t^{N-1}; q, t\big)}
\end{gather}
and call $P_{\lambda}(x_1, \dots, x_m; N, q, t)$ the \textit{Macdonald unitary character of rank $N$, number of variables $m$ and parametrized by $\lambda$}.
For simplicity of terminology, we call $P_{\lambda}(x_1, \dots, x_m; N, q, t)$ a \textit{Macdonald character} rather than a Macdonald unitary character. Observe that if $q, t \in \C$ are such that $|q|, |t|\in (0, 1)$, the evaluation identity for Macdonald polynomials, Theorem~\ref{evaluation}, shows that the denominator of~\eqref{normalmacdonald} is nonzero.
\end{df}

\section{Integral formulas for Macdonald characters of one variable}\label{macdonaldintegralsec}


In this section, assume $q$ is a real number in the interval $(0, 1)$. There will also be a parameter $\theta$, typically $\theta > 0$, but we also consider cases when $\theta$ is a complex number with $\Re\theta > 0$. In either case, the parameter $t = q^{\theta}$ satisf\/ies $|t| < 1$.


\subsection{Statements of the theorems}


The simplest contour integral representation is the following, which works only when $t = q^{\theta}$, $\theta\in\N$, and involves a closed contour around f\/initely many singularities.


\begin{thm}\label{macdonaldthm1}
Let $\theta\in\N$, $t = q^{\theta}$, $N\in\N$, $\lambda\in\GT_N$ and $x\in\C \setminus \big\{0, q, q^2, \dots, q^{\theta N - 1}\big\}$. Then
\begin{gather}
\frac{P_{\lambda}\big(x, t, t^2, \dots, t^{N-1}; q, t\big)}{P_{\lambda}\big(1, t, t^2, \dots, t^{N-1}; q, t\big)} \nonumber\\
\qquad {} = \ln(1/q)\prod_{i=1}^{\theta N - 1}{\frac{1 - q^i}{x - q^i}}\frac{1}{2\pi\sqrt{-1}}\oint_{\CC_0}{\frac{x^z}{\prod\limits_{i=1}^N\prod\limits_{j=0}^{\theta -1}{\big(1 - q^{z - (\lambda_i + \theta(N-i) + j)}\big)}}{\rm d}z},\label{macdonaldthm1eqn}
\end{gather}
where $\CC_0$ is a closed, positively oriented contour enclosing the real poles $\{\lambda_i + \theta(N-i) + j \colon i = 1, \dots, N, \, j = 0, \dots, \theta - 1\}$ of the integrand. For instance, the rectangular contour with vertices $-M-r\sqrt{-1}$, $ -M+r\sqrt{-1}$, $M+r\sqrt{-1}$ and $M-r\sqrt{-1}$, for any $-\frac{2\pi}{\ln{q}} > r > 0$ and any $M > \max\{0, -\lambda_N, \, \lambda_1 + \theta N - 1\}$, is a suitable contour.
\end{thm}

The following two theorems are analytic continuations, in the variable $\theta$, of Theorem~\ref{macdonaldthm1} above.

\begin{thm}\label{macdonaldthm2}
Let $\theta > 0$, $t = q^{\theta}$, $N\in\N$, $\lambda\in\GT_N$ and $x\in\C \setminus \{0\}$, $|x|\leq 1$. The integral below converges absolutely and the equality holds
\begin{gather}
\frac{P_{\lambda}\big(xt^{N-1}, t^{N-2}, \dots, t, 1; q, t\big)}{P_{\lambda}\big(t^{N-1}, t^{N-2}, \dots, t, 1; q, t\big)}\nonumber\\
 \qquad{} = \frac{\ln q}{1 - q}\frac{\big(xt^N; q\big)_{\infty}}{(xq; q)_{\infty}}\frac{\Gamma_q(\theta N)}{2\pi\sqrt{-1}}\int_{\CC^+}{x^z\prod_{i=1}^N{\frac{\Gamma_q(\lambda_i + \theta(N-i)-z)}{\Gamma_q(\lambda_i + \theta(N-i+1)-z)}}{\rm d}z}.\label{macdonaldthm2eqn}
\end{gather}
Contour $\CC^+$ is a positively oriented contour consisting of the segment $[M+r\sqrt{-1}, M-r\sqrt{-1}]$ and the horizontal lines $[M+r\sqrt{-1}, +\infty+r\sqrt{-1})$, $[M-r\sqrt{-1}, +\infty-r\sqrt{-1})$, for some $-\frac{\pi}{2\ln{q}} > r > 0$ and $\lambda_N > M$, see Fig.~{\rm \ref{fig:C}}. Observe that $\CC^+$ encloses all real poles of the integrand $($which accumulate at~$+\infty)$ and no other poles.
\end{thm}

The reader is referred to Appendix~\ref{app:qtheory} for a reminder of the def\/inition of the $q$-Gamma function, its zeroes and poles.
